% =====================================================================
%  c04 —— 植田一石《同调镜像对称》中译（第 9 节续文 + 第 10–11 节）
%  页范围：p016–p018
%  编号公式：15 条 —— (9.2) (10.1)–(10.14)
%  说明：BRIEF 原记 17 条，系 work/tags.txt 的计数含 p018 正文交叉引用
%        「(10.12) の右辺」并另有 1 条重复误检；按页面图像实际为 15 条。
%  节标题：p016 顶部为第 9 节续文（不写节标题）；第 10 节起于 p016、
%        第 11 节起于 p018，各写一次。
% =====================================================================

的函数 $f\colon\R^n\to\R^m$ 称为仿射函数（affine function）。这里，当 $A\in\operatorname{Mat}(m,n;\Z)$ 时，称 $f$ 为热带仿射函数（tropical affine function）；当 $A\in\operatorname{Mat}(m,n;\Z)$ 且 $b\in\Z^n$ 时，称 $f$ 为整仿射函数（integral affine function）\footnote{这里所说的热带仿射函数也常被称为整仿射函数。}。变换函数为仿射函数的流形称为仿射流形（affine manifold）。热带仿射流形与整仿射流形也类似地定义。切丛的联络称为仿射联络（affine connection），而仿射流形可刻画为具有无挠平坦仿射联络的流形。又，对热带仿射流形 $B^{\mathrm{sm}}$，其上的热带仿射函数所成的层 $\mathrm{Aff}_{B^{\mathrm{sm}}}$，成为把切丛的 $\Z$ 格 $T^{\Z}B^{\mathrm{sm}}\subset TB^{\mathrm{sm}}$ 视为局部常层所得的层、经常值层 $\R$ 作出的扩张；

\begin{equation}
    0\to\R\to\mathrm{Aff}_{B^{\mathrm{sm}}}\to T^{\Z}B^{\mathrm{sm}}\to 0. \tag{9.2}
\end{equation}

反之，若给定切丛的 $\Z$ 格 $T^{\Z}B^{\mathrm{sm}}\subset TB^{\mathrm{sm}}$ 及其经 $\R$ 的扩张 $\mathrm{Aff}_{B^{\mathrm{sm}}}$，则唯一确定使 $\mathrm{Aff}_{B^{\mathrm{sm}}}$ 成为热带仿射函数层的热带仿射结构。

给定完全可积系统时，作用坐标的任意性仅在于 $H_1(L;\Z)$ 的 $\Z$ 上有序基与 $U$ 的基点 $b_0$ 的更换。对作用坐标，前者以 $\mathrm{GL}_n(\Z)$、后者以 $\R^n$ 的平移作用，因此这一任意性由热带仿射变换群 $\mathrm{GL}_n(\Z)\ltimes\R^n$ 的作用给出。这说明用开集覆盖 $B^{\mathrm{sm}}$ 时，坐标变换成为热带仿射变换，所以 $B^{\mathrm{sm}}$ 上带有热带仿射结构。角坐标的任意性本质上是 $\pi|_{\pi^{-1}(U)}$ 的平凡化的任意性，这与截面的取法的自由度一致。

对热带仿射流形 $B^{\mathrm{sm}}$，整切丛 $T^{\Z}B^{\mathrm{sm}}\subset T$ 的对偶丛 $T^{*}_{\Z}B^{\mathrm{sm}}$ 被确定，作为其商，环面纤维丛 $X(B^{\mathrm{sm}}):=T^{*}B^{\mathrm{sm}}/T^{*}_{\Z}B^{\mathrm{sm}}$ 被确定。$X(B^{\mathrm{sm}})$ 与 $\pi^{-1}(B^{\mathrm{sm}})$ 能通过保持纤维的辛同胚相互转换的充要条件是：$\pi|_{\pi^{-1}(B^{\mathrm{sm}})}$ 具有截面。

当 $\pi|_{\pi^{-1}(B^{\mathrm{sm}})}$ 不具有截面时，$\pi^{-1}(B^{\mathrm{sm}})$ 是通过在各开集 $U_i$ 上取截面、并在 $U_i\cap U_j$ 上把它们粘合而得到的。借 $U_i$ 上的平凡化来看 $U_j$ 的截面，得到 $U_i\cap U_j$ 上的拉格朗日截面，它对应于 $U_i\cap U_j$ 上的闭 1 次微分形式。把这个闭 1 次微分形式记作 $f_{ij}$，则在 $U_i\cap U_j\cap U_k$ 上 $\alpha_{ijk}:=f_{ij}+f_{jk}-f_{ik}$ 成为热带仿射函数（微分后成为 $T^{*}_{\Z}(U_i\cap U_j\cap U_k)$ 的常数截面）。$\{\alpha_{ijk}\}_{i,j,k}$ 成为 \v{C}ech 上闭链（cocycle）。又，$f_{ij}$ 有加上常数的任意性，这成为上边缘（coboundary）。因此，若用 $\mathrm{Aff}$ 表示 $M$ 上的热带仿射函数所成的层，则 $B^{\mathrm{sm}}$ 上的拉格朗日环面纤维丛在保持纤维的辛同胚下的等价类，由层系数上同调 $H^2(B^{\mathrm{sm}},\mathrm{Aff})$ 分类 [Dui80]。

% [Tutor Note: 上段原文作「M 上のトロピカルアファイン関数のなす層」（p16），按上下文似应为「B^sm 上」；此处照原文保留，请审校。]

\section{Floer 上同调与 Fukaya 范畴}

周期解的存在是动力系统的中心问题之一。对任意紧辛流形 $(M,\omega)$，若把遍历 $M$ 上全体 $C^\infty$ 函数的临界点个数的下界记为 $c_M:=\inf_{f\in C^\infty(M)}\#\,\operatorname{Crit}(f)$，则 Arnold 曾猜想：$M$ 上任一周期 Hamilton 系统至少有 $c_M$ 个不动点。作为该猜想的推广，有 Arnold–Givental 猜想（Arnold–Givental conjecture），它断言：

% [Tutor Note: preamble 未定义「予想（conjecture）」环境；此处以 \para 手动标注原文编号「猜想 10.1」，以保持与原文编号一致，请 Lead 视需要增设 conjecture 环境。]
\para{猜想 10.1.}设 $(M,\omega)$ 为紧辛流形，$L$ 为 $M$ 的紧拉格朗日子流形。对依赖于时间的 Hamilton 函数 $H\colon M\times[0,1]\to\R$，记其生成的流为 $\psi\colon M\times[0,1]\to M$。假设由 $\psi_1:=\psi(-,t)$ 给出的 $L$ 的像 $\psi_1(L)$ 与 $L$ 横截相交，则 $\#(L\cap\psi_1(L))\ge\rank H^*(L;\Z/2\Z)$。

% [Tutor Note: 原文 p16 作「ψ₁ := ψ(−, t)」，下标 1 与 t 不一致；按文意或为 ψ_t := ψ(−, t)（或 ψ₁ := ψ(−, 1) 之误）。此处照原文保留。]

将 Arnold–Givental 猜想应用于 $(M\times M,\pi_1^*\omega_M-\pi_2^*\omega_M)$ 的对角集与 $(\id,\psi_1)$，即可推出弱 Arnold 猜想（若紧辛流形 $M$ 上周期 Hamilton 系统的所有不动点都非退化，则不动点的个数不少于 $\rank H^*(M;\Z/2\Z)$）。

Lagrange 相交 Floer 上同调，是在道路空间 $\Omega(L_0,L_1):=\{\ell\colon[0,1]\to M\mid\ell(0)\in L_0,\ \ell(1)\in L_1\}$ 的不变覆叠 $\Omega(L_0,L_1)^{\sim}$ 上、以作用函数
\begin{equation}
    \cA(\ell):=\int_{[0,1]\times[0,1]}\gamma^*\omega \tag{10.1}
\end{equation}
为 Morse 函数的 Morse 理论。

% [Tutor Note: 原文 p16「不変被覆」或为「普遍被覆（universal cover）」之误；此处照原文译作「不变覆叠」，请审校。]

这里，对 $\Omega(L_0,L_1)$ 的每个连通分支取定一个基点 $\ell_0$，并对 $\ell\in\Omega(L_0,L_1)$ 取从 $\ell_0$ 到 $\ell$ 的道路 $\gamma\colon[0,1]\to\Omega(L_0,L_1)$。由于 $\omega$ 是闭形式，由 Stokes 定理，$\cA$ 不依赖于 $\gamma$ 的取法，而定出 $\Omega(L_0,L_1)^{\sim}$ 上的函数。

在 $(M,\omega)$ 上取相容的近复结构，它就定出 $M$ 的 Riemann 度量 $g_M$。对 $\ell\in\Omega(L_0,L_1)$，有 $T_\ell\Omega(L_0,L_1)=\Gamma(\ell^*TM;T_{\ell(0)}L_0,T_{\ell(1)}L_1)$，对其元素 $V,W$，若定义
\begin{equation}
    g_{\Omega(L_0,L_1)}(V,W):=\int_0^1 g_M(V(t),W(t))\,dt \tag{10.2}
\end{equation}
则这给出 $\Omega(L_0,L_1)$ 的度量。

Gromov 的重要观察是：$\cA$ 的梯度流的轨迹成为近全纯映射。事实上，对 $\ell\in\Omega(L_0,L_1)$ 与 $V\in T_\ell\Omega(L_0,L_1)\cong\Gamma(\ell^*TM;T_{\ell(0)}L_0,T_{\ell(1)}L_1)$，令 $\ell_s=\exp_{\ell(t)}(sV(t))$，则
\begin{align}
    d\cA(V) &= \int_0^1 \omega\left(V,\frac{\partial\ell}{\partial t}\right)dt \tag{10.3}\\
    &= -\int_0^1 \omega\left(\frac{\partial\ell}{\partial t},V\right)dt \tag{10.4}\\
    &= -\int_0^1 g_M\left(J_M\left(\frac{\partial\ell}{\partial t}\right),V\right)dt \tag{10.5}\\
    &= g_{\Omega(L_0,L_1)}\left(-J_M\left(\frac{\partial\ell}{\partial t}\right),V\right) \tag{10.6}
\end{align}
即
\begin{equation}
    (\operatorname{grad}\cA)_\ell=-J_M\left(\frac{\partial\ell}{\partial t}\right) \tag{10.7}
\end{equation}
成立。因此，$\gamma\colon\R\to\Omega(L_0,L_1)$ 为 $\operatorname{grad}\cA$ 的积分曲线的充要条件是
\begin{equation}
    \frac{\partial\gamma}{\partial s}=-J_M\left(\frac{\partial\ell}{\partial t}\right) \tag{10.8}
\end{equation}
成立。把 $\gamma$ 视为从 $\Sigma:=\{t+\sqrt{-1}s\in\C\mid s\in[0,1],\ t\in\R\}$ 到 $M$ 的映射，将其改写，则
\begin{equation}
    -\gamma_*\circ J_\Sigma\left(\frac{\partial}{\partial t}\right)=-J_M\circ\gamma_*\left(\frac{\partial}{\partial t}\right) \tag{10.9}
\end{equation}
即
\begin{equation}
    \gamma_*\circ J_\Sigma=J_M\circ\gamma_* \tag{10.10}
\end{equation}
成立。

$\cA$ 的临界点是到 $L_0$ 与 $L_1$ 的交集 $L_0\cap L_1$ 的常值函数。若紧拉格朗日子流形 $L_0$ 与 $L_1$ 横截相交，则交集 $L_0\cap L_1$ 由有限个点组成。定义 $\Z/2\Z$ 上的 Novikov 域（Novikov field）为
\begin{equation}
    \Lambda=\left\{\sum_{i=0}^{\infty}a_iT^{\lambda_i}\ \middle|\ a_i\in\Z/2\Z,\ \lambda_i\in\R,\ \lim_{i\to\infty}\lambda_i\to\infty\right\} \tag{10.11}
\end{equation}
并对以 $L_0\cap L_1$ 为基的 $\Lambda$ 上的自由模赋予
\begin{equation}
    \mathbf{m}_1(p)=\sum_{q\in L_0\cap L_1}\sum_{[\varphi]\in\overline{\cM}_2(L_0,L_1;p,q)}T^{\langle[\varphi],\omega\rangle}q \tag{10.12}
\end{equation}
作为微分，所得之物记作 $\mathrm{CF}^*(L_0,L_1)$，称为 Lagrange 相交 Floer 复形（Lagrangian intersection Floer complex）。这里 $T$ 是 Novikov 域的不定元，$\overline{\cM}_2(L_0,L_1;p,q)$ 是从 $\Sigma$ 到 $M$ 的近全纯映射的模空间的稳定紧致化。$\langle[\varphi],\omega\rangle$ 是 $[\varphi]$ 所代表的 $\pi_2(M,L_0\cap L_1)$ 的元素与 $\omega\in H^2(M,L_0\cap L_1)$ 的配对，给出近全纯映射的能量（energy）。$\overline{\cM}_2(L_0,L_1;p,q)$ 的虚拟维数（virtual dimension）由指数定理计算，而求和只对虚拟维数为 0 的模空间的连通分支进行。此外，必要时对方程施加适当的扰动，使虚拟维数与真实维数一致。由 Gromov 紧性定理，能量上界被指定的模空间是紧的，因此 (10.12) 的右边作为 Novikov 域的元素是良定义的。复形的分次由 Lagrange 相交的 Maslov 指数（Maslov index）给出，一般取值于 $\Z/2\Z$。要将其提升到 $\Z$，需要在辛流形与拉格朗日子流形两者上附加称为分次（grading）的额外结构。

在适当条件下 $\mathbf{m}_1\circ\mathbf{m}_1=0$ 成立，其上同调 $\mathrm{HF}^*(L_0,L_1):=\operatorname{Ker}\mathbf{m}_1/\operatorname{Im}\mathbf{m}_1$ 称为 Lagrange 相交 Floer 上同调（Lagrangian intersection Floer cohomology）。Lagrange 相交 Floer 上同调对 Hamilton 同胚 $\phi\colon M\to M$ 具有不变性
\begin{equation}
    \mathrm{HF}(L_1,L_2)\cong\mathrm{HF}(L_1,\phi(L_2)) \tag{10.13}
\end{equation}
并在 $\phi(L)$ 与 $L$ 横截相交时再现 $L$ 的 $\Lambda$ 系数通常上同调（ordinary cohomology）；
\begin{equation}
    \mathrm{HF}^*(L,\phi(L);\Lambda)\cong H^*(L;\Lambda). \tag{10.14}
\end{equation}
Arnold–Givental 猜想由这些性质立即推出。

深谷 [Fuk93] 对给定的辛流形 $M$，提议考虑以 $M$ 的拉格朗日子流形为对象、以 Lagrange 相交 Floer 复形为态射空间的 $A_\infty$ 范畴\footnote{在 Fukaya 范畴中，只考虑 $\mathbf{m}_1$ 的上同调中的积 $[\mathbf{m}_2]$ 而忽略高阶积的东西，有时也称为 Donaldson–Fukaya 范畴（Donaldson–Fukaya category）。}。该范畴中的高阶积由从带标记圆盘出发的近全纯映射的计数来定义。这就是 Fukaya 范畴（Fukaya category），并被 Kontsevich [Kon95b] 立即用于同调镜像对称的表述\footnote{引入 Fukaya 范畴的最初动机之一，是把出现在带边 4 维流形的 Donaldson 理论中的 3 维流形的瞬子（instanton）Floer 上同调推广到带边 3 维流形。这与 Atiyah–Floer 猜想密切相关。}。

\section{Seidel 纲领}

要计算紧辛流形的 Fukaya 范畴，必须对以拉格朗日子流形为边界的近全纯圆盘进行计数，但从椭圆曲线情形中出现 $\vartheta$ 函数也可看出，这种计数往往是无限和，具体执行极为困难。

使射影流形的 Fukaya 范畴的处理成为可能的，是 Seidel 纲领，它由以下两个步骤构成：

\begin{itemize}[leftmargin=2.4em]
    \item 射影流形的 Fukaya 范畴是仿射流形的 Fukaya 范畴的形变。
    \item 仿射流形的 Fukaya 范畴可通过 Picard–Lefschetz 理论组合地计算。
\end{itemize}

Seidel 考虑：当辛流形 $X$ 是射影流形时，从 $X$ 中除去超平面截面 $D$ 而使其成为仿射流形。仿射流形是正合辛流形，与紧辛流形相比被期待更易处理。在射影流形 $X$ 中，超平面截面 $D$ 是辛形式 $\omega$ 的 Poincaré 对偶，因此 (10.12) 中 $T$ 的指数 $\langle[\varphi],\omega\rangle$ 成为 $[\varphi]$ 与 $D$ 的交点数。于是，Floer 复形定义在比 Novikov 域更小的收敛幂级数环 $(\Z/2\Z)[\![T]\!]$ 上，并且通过把 $0$ 代入 $T$ 可以再现 $X\setminus D$ 中的 Floer 复形。这表明 $X$ 的 Fukaya 范畴是 $X\setminus D$ 的 Fukaya 范畴的形变，开辟了用形变理论研究 Fukaya 范畴的道路 [Sei02]。

作为辛几何中的强有力工具，有 Lefschetz 纤维化（Lefschetz fibration）的概念。Lefschetz 纤维化是 Morse 函数在复几何中的类似物，是允许非退化\footnote{若适当取坐标后可表为 $\C^d\to\C$，$(z_1,\dots,z_d)\mapsto z_1^2+\cdots+z_d^2$ 的临界点，称为非退化（non-degenerate）。}临界点的纤维丛的推广概念。根据 Donaldson 的著名工作 [Don96]，已知任意辛流形适当爆破后具有 Lefschetz 纤维化结构。Seidel 引入了对全空间为正合辛流形的 Lefschetz 纤维化 $\pi\colon M\to S$ 的 Fukaya 范畴的变种 $\Fuk(\pi)$，

% CHUNK c04 DONE pages p016–p018
